\documentclass[aspectratio=169]{beamer}
\usepackage{kotex}
\usepackage{listings}
\usepackage[lighttt]{lmodern}
\usepackage{graphics}
\usepackage{amsmath}
\usepackage{mathtools}
\usepackage{xcolor}

\setsanshangulfont[ItalicFont={*},ItalicFeatures={FakeSlant=.22}]{Noto Sans CJK KR}
\title{All You Need Is Superword-Level Parallelism: \\ Systematic Control-Flow Vectorization with SLP (PLDI '22)}
\date{January 16, 2025}
\author[shortname]{Yishen Chen\inst{1} \and Charith Mendis\inst{2} \and Saman Amarasinghe\inst{1}}
\institute[shortinst]{\inst{1} Massachusetts Institute of Technology \and \inst{2} University of Illinois at Urbana-Champaign}

\usetheme{default}
\setbeamertemplate{itemize items}[circle]

\lstset{
  numbers=none, frame=single, showspaces=false,
  showstringspaces=false, showtabs=false, breaklines=true, showlines=true,
  breakatwhitespace=true, basicstyle=\scriptsize\ttfamily\color{black},
  keywordstyle=\bfseries, basewidth=0.5em,
  moredelim=**[is][\color{lightgray}]{(*@}{@*)},
  moredelim=**[is][\color{red}]{(|@}{@|)},
}

\begin{document}
\begin{frame}
  \titlepage
\end{frame}

\begin{frame}
  \frametitle{Outline}
  \tableofcontents
\end{frame}

\AtBeginSection[]
{
\begin{frame}
  \frametitle{Outline}
  \tableofcontents[currentsection]
\end{frame}
}

\section{Introduction and Motivation}
\subsection{Current State of Autovectorization}
\begin{frame}
  \frametitle{Current State of Autovectorization}
  \framesubtitle{A tale of two methodologies}

  \begin{itemize}
    \item Two methods of autovectorization: {\bfseries loop vectorization} and {\bfseries SLP vectorization}
    \pause
    \item Loop vectorization: targets {\bfseries loop-level parallelism}
    \begin{itemize}
      \item Influenced by long-vector architectures (e.g., CRAY)
      \item Works by widening scalar instructions in loop body to vectorized counterparts
      \item Requires complex dependence analyses
    \end{itemize}
    \pause
    \item Superword-level parallelism vectorization: exploits {\bfseries short-vector parallelism}
    \begin{itemize}
      \item Influenced by modern vector extensions (e.g., Intel AVX)
      \item Works by packing independent and isomorphic groups of instructions
      \item More flexible, but harder to handle control flow
    \end{itemize}
    \pause
    \item The two techniques have their advantages and disadvantages
    \begin{itemize}
      \item Production compilers use both techniques in separate passes
    \end{itemize}
  \end{itemize}
\end{frame}

\begin{frame}[fragile]
  \frametitle{Current State of Autovectorization}
  \framesubtitle{Loop vectorization in action}

  \begin{itemize}
    \item SAXPY loop widened by factor of 8
  \end{itemize}
  \begin{figure}
    \centering
    \begin{minipage}{.45\textwidth}
      \begin{lstlisting}[language=C]
void saxpy(float a, float x[], float y[], int n) {
  for (int i = 0; i < n; i++)
    y[i] = a * x[i] + y[i];
}
      \end{lstlisting}
    \end{minipage}
    \hspace{0.2cm}
    \pause
    \begin{minipage}{.45\textwidth}
      \begin{lstlisting}[language=C]
void saxpy(float a, float x[], float y[], int n) {
  int i;
  for (i = 0; i + 8 <= n; i += 8)
    y[i:i + 8] = splat(a, 8) * x[i:i + 8] + y[i:i + 8];
  for (; i < n; i++)
    y[i] = a * x[i] + y[i];
}
      \end{lstlisting}
    \end{minipage}
  \end{figure}
\end{frame}

\begin{frame}[fragile]
  \frametitle{Current State of Autovectorization}
  \framesubtitle{SLP vectorization in action}

  \begin{itemize}
    \item Four elements grouped and processed together
  \end{itemize}
  \begin{figure}
    \centering
    \begin{minipage}{.45\textwidth}
      \begin{lstlisting}[language=C]
void foo(int a1, int a2, int b1, int b2, int *A) {
  A[0] = a1 * (a1 + b1);
  A[1] = a2 * (a2 + b2);
  A[2] = a1 * (a1 + b1);
  A[3] = a2 * (a2 + b2);
}
      \end{lstlisting}
    \end{minipage}
    \hspace{0.2cm}
    \pause
    \begin{minipage}{.45\textwidth}
      \begin{lstlisting}[language=C]
void foo(int a1, int a2, int b1, int b2, int *A) {
  va = {a1, a2, a1, a2};
  vb = {b1, b2, b1, b2};
  A[0:4] = va * (va + vb);
}
      \end{lstlisting}
    \end{minipage}
  \end{figure}
\end{frame}

\subsection{Flexibility of SLP Vectorization}
\begin{frame}[fragile]
  \frametitle{Flexibility of SLP Vectorization}
  \framesubtitle{Inefficiency of loop vectorization}

  \begin{itemize}
    \item Naive vectorization of unrolled SAXPY code
  \end{itemize}
  \begin{figure}
    \centering
    \begin{minipage}{.45\textwidth}
      \begin{lstlisting}[language=C]
void saxpy(float a, float x[], float y[], int n) {
  for (int i = 0; i + 2 <= n; i += 2) {
    y[i] = a * x[i] + y[i];
    y[i + 1] = a * x[i + 1] + y[i + 1];
  }
  // handle leftover elements
}
      \end{lstlisting}
    \end{minipage}
    \hspace{0.2cm}
    \pause
    \begin{minipage}{.45\textwidth}
      \begin{lstlisting}[language=C]
void saxpy(float a, float x[], float y[], int n) {
  int i;
  for (i = 0; i + 16 <= n; i += 16) {
    vx = x[i:i + 16]; vy = y[i:i + 16];
    vx_even = shuffle_even(vx);
    vy_even = shuffle_even(vy);
    vx_odd = shuffle_odd(vx);
    vy_odd = shuffle_odd(vy);
    y[i:i + 16] = interleave(
      splat(a, 8) * vx_even + vy_even,
      splat(a, 8) * vx_odd + vy_odd
    );
  }
  // handle leftover elements
}
      \end{lstlisting}
    \end{minipage}
  \end{figure}
\end{frame}

\begin{frame}[fragile]
  \frametitle{Flexibility of SLP Vectorization}
  \framesubtitle{Flexibility of unroll-and-SLP}

  \begin{itemize}
    \item After unrolling 4 times, SLP vectorization leads to better code
  \end{itemize}
  \begin{figure}
    \centering
    \begin{minipage}{.45\textwidth}
      \begin{lstlisting}[language=C]
void saxpy(float a, float x[], float y[], int n) {
  for (int i = 0; i < n; i += 8) {
    y[i] = a * x[i] + y[i];
    y[i + 1] = a * x[i + 1] + y[i + 1];
    y[i + 2] = a * x[i + 2] + y[i + 2];
    y[i + 3] = a * x[i + 3] + y[i + 3];
    y[i + 4] = a * x[i + 4] + y[i + 4];
    y[i + 5] = a * x[i + 5] + y[i + 5];
    y[i + 6] = a * x[i + 6] + y[i + 6];
    y[i + 7] = a * x[i + 7] + y[i + 7];
  }
  // handle leftover elements
}
      \end{lstlisting}
    \end{minipage}
    \hspace{0.2cm}
    \pause
    \begin{minipage}{.45\textwidth}
      \begin{lstlisting}[language=C]
void saxpy(float a, float x[], float y[], int n) {
  int i;
  for (i = 0; i + 8 <= n; i += 8)
    y[i:i + 8] = splat(a, 8) * x[i:i + 8] + y[i:i + 8];
  // handle leftover elements
}
      \end{lstlisting}
    \end{minipage}
  \end{figure}
\end{frame}

\subsection{Key Contributions}
\begin{frame}
  \frametitle{TL;DR}
  \framesubtitle{Key contributions}

  \begin{enumerate}
    \item {\bfseries Predicated SSA} for simplifying code motion across arbitrary control flow regions
    \item {\bfseries SuperVectorization}, a generalized form of SLP vectorization capable of handling inner and outer loops as well as straight-line code
    \item Prototype implementation which outperforms LLVM's vectorization passes
  \end{enumerate}
\end{frame}

\section{Groundwork for SuperVectorization}
\subsection{Handling Control Flow in SLP Vectorization}
\begin{frame}[fragile]
  \frametitle{Shortcomings of Conventional SLP Vectorization}

  \begin{itemize}
    \item Conventional SLP vectorization has been limited to straight-line code
    \begin{itemize}
      \item However, SLP is not just about similar instructions in a single basic block!
    \end{itemize}
  \end{itemize}

  \pause
  \begin{lstlisting}[language=C]
void search(int haystack[], int needles[2], int needle_idxs[2], int n) {
  for (int i = 0; i < n; i++)
    if (haystack[i] == needles[0]) {
      needle_idxs[0] = i;
      break;
    }
  for (int i = 0; i < n; i++)
    if (haystack[i] == needles[1]) {
      needle_idxs[1] = i;
      break;
    }
}
  \end{lstlisting}
  \begin{itemize}
    \item The two loops are independent and, thus can be run in parallel
  \end{itemize}
\end{frame}

\begin{frame}
  \frametitle{Handling Control Flow in SLP Vectorization}
  \framesubtitle{Key challenges}

  \begin{itemize}
    \item {\bfseries Code motion} across basic blocks is challenging
    \begin{itemize}
      \item Packing instructions involve code motion
      \item Exploiting SLP across basic block boundaries means code motion across BBs
      \item In traditional IR with CFG this requires restructuring CFG, which is complex
    \end{itemize}
    \item {\bfseries If-conversion}: partially convert control flow into straight-line code
    \begin{itemize}
      \item \(\phi\)-nodes are converted into \texttt{select} instructions
      \item Analogous to converting \texttt{if (cond) f(); else g();} into \texttt{cond ? f() : g();}
      \item All branches have to be executed regardless of condition
    \end{itemize}
    \pause
    \item Solution: {\bfseries Predicated SSA} for code motion and conditional if-conversion
  \end{itemize}
\end{frame}

\subsection{Introduction to Predicated SSA}
\begin{frame}
  \frametitle{Introduction to Predicated SSA}
  \framesubtitle{Definition}

  \begin{itemize}
    \item Basic idea: ``flatten'' the CFG to expose SLP, while retaining control dependences
    \pause
    \item Remove distinctions of basic blocks by introducing {\bfseries control predicates}
    \begin{itemize}
      \item Annotates each instruction as boolean flags determining whether to execute
    \end{itemize}
    \pause
    \item {\bfseries Differentiate \(\phi\)-nodes} to capture its precise semantics
    \begin{itemize}
      \item Two main uses of \(\phi\)-nodes: loop headers and control-flow joins
      \item \(\mu\)-nodes for uses in loop header, gated \(\phi\)-nodes for control-flow joins
    \end{itemize}
  \end{itemize}
\end{frame}

\begin{frame}
  \frametitle{Introduction to Predicated SSA}
  \framesubtitle{Definition}

  \begin{align*}
    \mathsf{fn} &\Coloneqq \mathsf{item}_1 : \mathsf{p}_1, \dots, \mathsf{item}_n : \mathsf{p}_n \\
    \mathsf{c} &\Coloneqq \textit{branch conditions} \\
    \mathsf{p} &\Coloneqq \textbf{true}\; |\; \mathsf{c}\; |\; \bar{\mathsf{c}}\; |\; \mathsf{p}_1 \land \mathsf{p}_2\; |\; \mathsf{p}_1 \lor \mathsf{p}_2 \\
    \mathsf{item} &\Coloneqq \textit{instruction}\; |\; \mathsf{loop} \\
    \mathsf{loop} &\Coloneqq \textbf{with}\; \mathsf{v}_1 = \mathsf{\mu}_1, \dots \mathsf{v}_m = \mathsf{\mu}_m\; \textbf{do}\; \mathsf{item}_1 : \mathsf{p}_1, \dots, \mathsf{item}_n : \mathsf{p}_n\; \textbf{while}\; \mathsf{p}_\text{cont} \\
    \mathsf{\mu} &\Coloneqq \textbf{mu}(\mathsf{v}_\text{init}, \mathsf{v}_\text{rec}) \\
    \mathsf{\phi} &\Coloneqq \textbf{phi}(\mathsf{v}_1 : \mathsf{p}_1, \dots, \mathsf{v}_n : \mathsf{p}_n) \\
    \mathsf{v} &\Coloneqq \mathsf{\mu}\; |\; \textit{instruction}\; |\; \textit{constant}\; |\; \textit{argument}
  \end{align*}
\end{frame}

\begin{frame}[fragile]
  \frametitle{Introduction to Predicated SSA}
  \framesubtitle{Example}

  \begin{figure}
    \centering
    \begin{minipage}{.45\textwidth}
      \begin{lstlisting}[language=C]
for (int i = 0; i < n; i++)
  if (haystack[i] == needles[0]) {
    needle_idxs[0] = i;
    break;
  }
      \end{lstlisting}
    \end{minipage}
    \hspace{0.2cm}
    \begin{minipage}{.45\textwidth}
      \begin{lstlisting}
with i = mu(0, i') do
  t = load haystack[i]      : true
  needle = load needles[0]  : true
  found = cmp eq, t, needle : true
  i' = add i, 1             : true
  not_found = not found     : true
  lt_n = cmp lt i', n       : true
while not_found and lt_n    : true
store i, needle_idxs[0]     : found
      \end{lstlisting}
    \end{minipage}
  \end{figure}
\end{frame}

\section{SuperVectorization}
\begin{frame}
  \frametitle{Vectorization Workflow}
  \framesubtitle{SuperVectorization at a glance}

  \begin{enumerate}
    \item Converting CFG to Predicated SSA
    \item Unroll and Packing
    \item Loop transformations
    \item Vector instruction generation
    \item Converting Predicated SSA to CFG
  \end{enumerate}
\end{frame}

\subsection{Conversion to Predicated SSA}
\begin{frame}
  \frametitle{Conversion to Predicated SSA}
  \framesubtitle{High-level overview}

  \begin{enumerate}
    \item Convert input IR into a canonical form resembling \texttt{do-while} loops
    \begin{itemize}
      \item Single incoming edge and back edge, dedicated loop header, preheader, and latch
    \end{itemize}
    \pause
    \item Classify \(\phi\)-nodes according to its uses
    \begin{itemize}
      \item Forward \(\phi\)-nodes: use gated \(\phi\)-nodes
      \item Loop header \(\phi\)-nodes: use \(\mu\)-nodes
    \end{itemize}
    \pause
    \item Assign and calculate control predicates for all blocks
    \begin{itemize}
      \item Disjunction of branch conditions of control-dependent blocks' post-dominant successors
      \item All instructions in a block share one control predicate
    \end{itemize}
  \end{enumerate}
\end{frame}

\begin{frame}
  \frametitle{Conversion to Predicated SSA}
  \framesubtitle{Details of control predicate calculation}

  \begin{itemize}
    \item For a basic block \(b\) to execute, two conditions must be satisfied:
    \begin{enumerate}
      \item Control flow reaches some block \(b'\), a control-dependent basic block of \(b\)
      \item \(b'\) takes a branch that leads to \(b\)
    \end{enumerate}
    \pause
    \item Condition 1: check using {\bfseries post-dominance frontier}
    \pause
    \item Condition 2: simple to consider when all basic blocks have at most two successors
    \begin{itemize}
      \item i.e., no \texttt{switch} and indirect branches
      \pause
      \item Recall that \(b'\) is in the post-dominance frontier of \(b\)
      \item There exists an edge \(b' \to b''\) where \(b''\) is a successor of \(b'\) and post-dominant on \(b\)
      \item Such \(b''\) is unique as \(b'\) has at most two outgoing edges
    \end{itemize}
  \end{itemize}
\end{frame}

\begin{frame}[fragile]
  \frametitle{Conversion to Predicated SSA}
  \framesubtitle{Example}

  \begin{figure}
    \centering
    \begin{minipage}{0.45\textwidth}
      \begin{lstlisting}[language=C]
void search(int haystack[], int needles[2], int needle_idxs[2], int n) {
  for (int j = 0; j < n; j++)
    if (haystack[j] != needles[0]) {
      needle_idxs[0] = j;
      break;
    }
}
      \end{lstlisting}
    \end{minipage}
    \hspace{0.2cm}
    \begin{minipage}{0.45\textwidth}
      \includegraphics<1>[width=1\textwidth]{search_1x_cfg.pdf}%
      \includegraphics<2>[width=1\textwidth]{search_1x_cfg.if.then.pdf}%
      \includegraphics<3>[width=1\textwidth]{search_1x_cfg.if.then.postdom.pdf}%
    \end{minipage}
  \end{figure}
\end{frame}

\subsection{Performing Vectorization}
\begin{frame}
  \frametitle{Vector Packing}
  \framesubtitle{Overview}

  \begin{itemize}
    \item Capture SLP revealed by conversion to Predicated SSA and unrolling
    \item SuperVectorization is compatible with arbitrary packing heuristics
    \pause
    \begin{itemize}
      \item The authors adapt the bottom-up SLP heuristic, which is simple and commonly used
      \item Starting from \textit{seeds} (i.e., contiguous stores, reductions), traverse use-def chains
      \item Stop when nonvectorizable instructions are found, check dependence cycles and cost
    \end{itemize}
    \pause
    \item The packs are lowered into instructions using {\bfseries loop transforms}, {\bfseries scheduling}, and {\bfseries instruction selection}
  \end{itemize}
\end{frame}

\begin{frame}[fragile]
  \frametitle{Vector Packing}
  \framesubtitle{Example}

  \begin{figure}
    \centering
    \begin{minipage}{0.45\textwidth}
      \begin{lstlisting}[language=C]
void search(int haystack[], int needles[2], int needle_idxs[2], int n) {
  for (int j = 0; j < n; j++)
    if (haystack[j] != needles[0]) {
      needle_idxs[0] = j;
      break;
    }
  for (int j = 0; j < n; j++)
    if (haystack[j] != needles[1]) {
      needle_idxs[1] = j;
      break;
    }
}
      \end{lstlisting}
    \end{minipage}
    \hspace{0.2cm}
    \pause
    \begin{onlyenv}<1-2>
      \begin{minipage}{0.45\textwidth}
        \begin{lstlisting}
with i = mu(0, i') do
  t = load haystack[i]      : true
  needle = load needles[0]  : true
  found = cmp eq, t, needle : true
  i' = add i, 1             : true
  not_found = not found     : true
  lt_n = cmp lt i', n       : true
while not_found and lt_n    : true
store i, needle_idxs[0]     : found
with i2 = mu(0, i2') do
  t2 = load haystack[i2]       : true
  needle2 = load needles[1]    : true
  found2 = cmp eq, t2, needle2 : true
  i2' = add i2, 1              : true
  not_found2 = not found2      : true
  lt_n2 = cmp lt i2', n        : true
while not_found2 and lt_n2     : true
store i2, needle_idxs[1]       : found2
        \end{lstlisting}
      \end{minipage}
    \end{onlyenv}%
    \begin{onlyenv}<3>
      \begin{minipage}{0.45\textwidth}
        \begin{lstlisting}
(*@with i = mu(0, i') do@*)
(*@  t = load haystack[i]      : true@*)
(*@  needle = load needles[0]  : true@*)
(*@  found = cmp eq, t, needle : true@*)
(*@  i' = add i, 1             : true@*)
(*@  not_found = not found     : true@*)
(*@  lt_n = cmp lt i', n       : true@*)
(*@while not_found and lt_n    : true@*)
(|@store i, needle_idxs[0]     : found@|)
(*@with i2 = mu(0, i2') do@*)
(*@  t2 = load haystack[i2]       : true@*)
(*@  needle2 = load needles[1]    : true@*)
(*@  found2 = cmp eq, t2, needle2 : true@*)
(*@  i2' = add i2, 1              : true@*)
(*@  not_found2 = not found2      : true@*)
(*@  lt_n2 = cmp lt i2', n        : true@*)
(*@while not_found2 and lt_n2     : true@*)
(|@store i2, needle_idxs[1]       : found2@|)
        \end{lstlisting}
      \end{minipage}
    \end{onlyenv}%
    \begin{onlyenv}<4>
      \begin{minipage}{0.45\textwidth}
        \begin{lstlisting}
(*@with@*) i = mu(0, i') (*@do@*)
(*@  t = load haystack[i]      : true@*)
(*@  needle = load needles[0]  : true@*)
(*@  found = cmp eq, t, needle : true@*)
(*@  i' = add i, 1             : true@*)
(*@  not_found = not found     : true@*)
(*@  lt_n = cmp lt i', n       : true@*)
(*@while not_found and lt_n    : true@*)
(|@store i, needle_idxs[0]     : found@|)
(*@with@*) i2 = mu(0, i2') (*@do@*)
(*@  t2 = load haystack[i2]       : true@*)
(*@  needle2 = load needles[1]    : true@*)
(*@  found2 = cmp eq, t2, needle2 : true@*)
(*@  i2' = add i2, 1              : true@*)
(*@  not_found2 = not found2      : true@*)
(*@  lt_n2 = cmp lt i2', n        : true@*)
(*@while not_found2 and lt_n2     : true@*)
(|@store i2, needle_idxs[1]       : found2@|)
        \end{lstlisting}
      \end{minipage}
    \end{onlyenv}%
    \begin{onlyenv}<5>
      \begin{minipage}{0.45\textwidth}
        \begin{lstlisting}
(*@with@*) i = mu(0, i') (*@do@*)
(*@  t = load haystack[i]      : true@*)
(*@  needle = load needles[0]  : true@*)
(*@  found = cmp eq, t, needle : true@*)
  i' = add i, 1             : true
(*@  not_found = not found     : true@*)
(*@  lt_n = cmp lt i', n       : true@*)
(*@while not_found and lt_n    : true@*)
(|@store i, needle_idxs[0]     : found@|)
(*@with@*) i2 = mu(0, i2') (*@do@*)
(*@  t2 = load haystack[i2]       : true@*)
(*@  needle2 = load needles[1]    : true@*)
(*@  found2 = cmp eq, t2, needle2 : true@*)
  i2' = add i2, 1              : true
(*@  not_found2 = not found2      : true@*)
(*@  lt_n2 = cmp lt i2', n        : true@*)
(*@while not_found2 and lt_n2     : true@*)
(|@store i2, needle_idxs[1]       : found2@|)
        \end{lstlisting}
      \end{minipage}
    \end{onlyenv}%
    \begin{onlyenv}<6>
      \begin{minipage}{0.45\textwidth}
        \begin{lstlisting}
(*@with@*) i = mu(0, i') (*@do@*)
(*@  t = load haystack[i]      : true@*)
(*@  needle = load needles[0]  : true@*)
(*@  found = cmp eq, t, needle : true@*)
  i' = add i, 1             : true
(*@  not_found = not found     : true@*)
  lt_n = cmp lt i', n       : true
(*@while not_found and lt_n    : true@*)
(|@store i, needle_idxs[0]     : found@|)
(*@with@*) i2 = mu(0, i2') (*@do@*)
(*@  t2 = load haystack[i2]       : true@*)
(*@  needle2 = load needles[1]    : true@*)
(*@  found2 = cmp eq, t2, needle2 : true@*)
  i2' = add i2, 1              : true
(*@  not_found2 = not found2      : true@*)
  lt_n2 = cmp lt i2', n        : true
(*@while not_found2 and lt_n2     : true@*)
(|@store i2, needle_idxs[1]       : found2@|)
        \end{lstlisting}
      \end{minipage}
    \end{onlyenv}%
  \end{figure}
\end{frame}

\begin{frame}
  \frametitle{Lowering Packs into Instructions}
  \framesubtitle{Loop transforms: fusion and co-iteration}

  \begin{itemize}
    \item Predicate SSA allows packing from different loops
    \begin{itemize}
      \item However, the resulting vector instructions must be in the same loop
    \end{itemize}
    \pause
    \item Two ways of loop transformation: {\bfseries fusion} and {\bfseries co-iteration}
    \begin{itemize}
      \item Apply fusion if loops are independent, have the same trip-count/condition, and parent loops can be fused
      \item Co-iterate loops if they do not meet the conditions of fusion
    \end{itemize}
    \pause
    \item Co-iteration: run the loop until all conditions of the original loops become false
    \begin{itemize}
      \item Guard the loop bodies' execution using strengthened control predicates
    \end{itemize}
  \end{itemize}
\end{frame}

\begin{frame}
  \frametitle{Lowering Packs into Instructions}
  \framesubtitle{Scheduling and instruction selection}

  \begin{itemize}
    \item {\bfseries Scheduling}: reorder instructions to satisfy dependences introduced by packing
    \begin{itemize}
      \item Create a dependency graph of instructions and loops
      \item Since cyclic dependences have been checked before packing, the graph is DAG
      \item Topological sort on the dependency graph to obtain dependency-preserving order
    \end{itemize}
    \pause
    \item Special handling for load/store and \(\phi\)-node packs in {\bfseries instruction selection}
    \begin{itemize}
      \item For other instructions, naive packing into vector registers is fine
      \item Use masked loads and stores where masks are set according to control predicates
      \item When lowering \(\phi\)-nodes, emit \texttt{select} instructions if necessary
    \end{itemize}
  \end{itemize}
\end{frame}

\subsection{Converting Back to CFG}
\begin{frame}
  \frametitle{Converting Back to CFG}

  \begin{itemize}
    \item After lowering the packs, the work on Predicated SSA is done
    \pause
    \item Features specific to Predicated SSA are lowered into CFG constructs
    \begin{enumerate}
      \item Replace gated \(\phi\)-nodes with move instructions, ignoring SSA invariant for now
      \item Reconstruct basic blocks by querying {\bfseries block builder} with control predicates
      \item Generate back edges and loop exits
    \end{enumerate}
    \pause
    \item Block builder: a data structure for mapping control predicates into basic blocks
    \begin{itemize}
      \item When presented with control predicate \(p\), the following cases are possible:
    \end{itemize}
  \end{itemize}

  \begin{table}[]
  \begin{tabular}{|r|l|}
  \hline
  \(p\) & Block builder behavior \\ \hline
  \(\mathbf{true}\) & return top-level block \\
  \(p' \land c\) & add branch with condition \(c\) to \(p'\) block \\
  \(p_1 \lor \dots \lor p_n\) & join control-flow from \(p_1, \dots, p_n\) blocks \\ \hline
  \end{tabular}
  \end{table}
\end{frame}

\subsection{Performance Evaluation}
\begin{frame}
  \frametitle{Performance Evaluation}

  \begin{itemize}
    \item SuperVectorization outperformed LLVM's default \texttt{-O3} pipeline
    \begin{itemize}
      \item \({}^*\)Manually vectorized some functions
    \end{itemize}
  \end{itemize}
  \begin{figure}
  \centering
  \begin{tabular}{|r|ccc|}
  \hline
  Benchmark & TSVC & PolyBench & ISPC\({}^*\) \\ \hline
  Geomean speedup & \(1.04\times\) & \(1.46\times\) & \(1.88\times\) \\ \hline
  \end{tabular}
  \end{figure}
  \pause
  \begin{itemize}
    \item In some cases, performance regressed significantly (\(<0.2\times\))
    \begin{itemize}
      \item Missed vectorization opportunity due to data dependences
      \item Sub-optimal instruction selection (e.g., scalarized store instead of scatter)
      \item Incomplete handling of some reduction patterns
    \end{itemize}
  \end{itemize}
\end{frame}

\section{Takeaways}
\begin{frame}
  \frametitle{Takeaways}

  \begin{itemize}
    \item Dedicated representation facilitating code motion
    \begin{itemize}
      \item Thinking outside of CFG and basic blocks might reveal new opportunities
    \end{itemize}
    \pause
    \item Flexibility of SLP
    \begin{itemize}
      \item Exploiting SLP spanning multiple control flow regions had led to speedups
    \end{itemize}
    \pause
    \item Future directions
    \begin{itemize}
      \item Reduce dependence on loop unrolling for scalable vectorization support
    \end{itemize}
  \end{itemize}
\end{frame}
\end{document}
